% Generated by tools/edn_latex.py from reports/edn.md.
% Do not edit: change reports/edn.md and run `make edn`.
\ednentry{
  id = {EDN-31},
  title = {The processed-data contract is a hand-written \texorpdfstring{\texttt{schema.\allowbreak{}py}}{schema.py}, not Pandera},
  date = {2026-09-29},
  milestone = {M2: Reproducibility (shapes M3: Quality Assurance)},
  topic = {Data Validation, Testing Strategy},
  participants = {@lukas2510},
  decision = {\texttt{recommenditos/\allowbreak{}schema.\allowbreak{}py} holds the structural contract -- column names, dtypes and nullability for the raw, interim and processed frames -- as frozen dataclasses with a \texttt{validate} and a \texttt{conform} method, written by hand rather than with a validation library. Value rules stay in the Great Expectations suites. Every stage validates after reading and conforms before writing.},
  rationale = {\emph{Alternatives considered:}
\begin{itemize}[leftmargin=*, topsep=2pt]
\item \textbf{Option A (chosen): a hand-written \texttt{schema.\allowbreak{}py}.}
\newline Pros: no dependency, so the same module can later be imported by the API, which NFR-04 caps at 1 GB and which already has Pydantic for its own request schemas; we control the failure message, and it reports every problem at once with the offending row positions; and \texttt{conform} casts to the declared dtypes before each write, which is what stops Parquet's dtype drift between stages -- a library that only validates cannot do that. Great Expectations already gives us a declarative vocabulary for value rules, so a second declarative framework would overlap it.
\newline Cons: about 250 lines we own and test ourselves.
\item \textbf{Option B: Pandera.}
\newline Pros: battle-tested, declarative, less code to own, good lazy error reports as a tidy DataFrame.
\newline Cons: a third validation system next to Great Expectations and Pydantic; it would travel into the API image; and it does not check the datetime unit at all, so a \texttt{datetime64[ns]} declaration would be documentation rather than a guarantee -- exactly the failure this contract exists to prevent.
\item \textbf{Option C: Great Expectations alone, with no structural contract in code.}
\newline Pros: one tool, and Data Docs for free.
\newline Cons: the contract would exist only in a generated \texttt{gx/\allowbreak{}} store whose config carries a fresh UUID per object on every run, so it is awkward to diff and keep honest in git. It also arrives too late: the stage tickets need something to code against in week one, and Great Expectations is a later ticket.
\end{itemize}
\par
\emph{Why:} Structure and value are different questions. A structural break is a bug in our code and should fail the stage that caused it; a value break is a change in the data and belongs in an expectation suite with a \texttt{mostly=\allowbreak{}} tolerance. Splitting them that way let the contract land before Great Expectations does, which is what unblocks the parallel work. The decisive technical point is \texttt{conform}: Parquet preserves whatever dtype it is handed rather than normalising it, so without an explicit cast at each boundary the frames drift apart and the stage that notices is never the stage that caused it. Because \texttt{conform} selects only the contract's columns, the PII columns NFR-08 forbids are kept out structurally rather than by every stage remembering to drop them.},
  ai = {Information seeking, Alternative generation, Alternative assessment, Recommendation, Solution generation},
  response = {Accepted},
  assessment = {AI ran the compatibility experiments rather than arguing from documentation: it installed Pandera, Great Expectations and pyarrow against our pandas 3.0.6 and measured what survives a Parquet round trip, finding that an \texttt{object} string column silently becomes \texttt{str} across the boundary (so a schema declaring \texttt{object} passes upstream and fails downstream), that Pandera ignores the datetime unit entirely, and that the Great Expectations store regenerates with different UUIDs on every run. Those three measurements are what decided the option, and none of them was predictable from the documentation. Lukas was given the three options with these findings and chose A.},
  aievidence = {Claude Code session on 2026-09-29: AI ran the probes in a scratch environment, reported the before/after dtype table for the Parquet round trip and the sixteen declared-versus-actual datetime combinations Pandera accepts, and recommended the hand-written contract; Lukas chose it.},
  evidence = {\href{https://github.com/mlops-2627q1-mds-upc/MLOPS_Recommenditos/blob/main/recommenditos/schema.py}{\texttt{recommenditos/\allowbreak{}schema.\allowbreak{}py}}; \href{https://github.com/mlops-2627q1-mds-upc/MLOPS_Recommenditos/blob/main/tests/test_schema.py}{\texttt{tests/\allowbreak{}test\_\allowbreak{}schema.\allowbreak{}py}}; \href{https://github.com/mlops-2627q1-mds-upc/MLOPS_Recommenditos/issues/32}{issue \#32}.},
}
